% !PRINTED 237
\weberpage{237}

In this transformation the assumption has been tacitly made that the quantities denoted by $m$ and $c$ in (5) are not zero.

If $m = 0$, then in formula (4) $p\,u^2 + 2\,q\,u\,v + r\,v^2$ is divisible by $a\,u + b\,v$; hence, at least for some of the roots, the equation
\[
a\,u + b\,v = 0,
\]
that is, an equation of degree no higher than the fourth, must hold.

If $c = 0$, so that $p\,u^2 + 2\,q\,u\,v + r\,v^2$ is the square of a linear function, then it follows from (7) that, if we set
\[
y = a'\,u + b'\,v
\]
then
\[
S(y) = 0, \quad S(y^2) = 0, \quad S(y^3) = 0
\]
and hence that the \gesperrt{Bring-Jerrard} form can be produced in a rational way.

Finally, we may summarise the results of these considerations as follows:

\gesperrt{The principal equation of the fifth degree can be transformed, by a transformation which contains as its only irrationality the square root of the discriminant, into a normal form with one parameter.}

\begin{center}
\rule{0.25\linewidth}{0.4pt}
\end{center}
